\documentclass[11pt]{article}

\usepackage[a4paper,margin=1in]{geometry}
\usepackage[T1]{fontenc}
\usepackage[utf8]{inputenc}
\usepackage{amsmath}
\usepackage{graphicx}
\usepackage{booktabs}
\usepackage{tabularx}
\usepackage{array}
\usepackage{subcaption}
\usepackage{float}
\usepackage{hyperref}
\usepackage{enumitem}

\hypersetup{
    colorlinks=true,
    linkcolor=blue,
    citecolor=blue,
    urlcolor=blue
}

\graphicspath{{eps/}}

\title{Finite-Blocklength Latency Minimization for Uplink and Downlink Communications:\\Per-User Scheduling Results}
\author{}
\date{July 6, 2026}

\begin{document}

\maketitle

\begin{abstract}
This document focuses only on per-user scheduling behavior. Throughout, the
\emph{Training Only Method} is used as the benchmark for the
\emph{Training + Testing Method}. The clearest scheduling gap appears in the
downlink. In payload completion, the benchmark finishes with block counts
\((4,3,3)\), whereas the \emph{Training + Testing Method} expands that to
\((9,3,7)\). In the uplink, the scheduling structures are already close, and
the main visible difference is limited to a single extra payload block for
user 1 under the \emph{Training + Testing Method}.
\end{abstract}

\section{Scope and Naming}

This document is intentionally schedule-only. The terminology is aligned with
the companion theory report:
\begin{itemize}[leftmargin=1.5em]
    \item \textbf{Training Only Method}: direct online optimization during scheduling,
    \item \textbf{Training + Testing Method}: a separate training stage followed by testing-stage inference,
    \item \textbf{Payload completion}: the scenario where each user continues until its remaining payload is exhausted.
\end{itemize}

The \emph{Training Only Method} is treated as the benchmark, and the
\emph{Training + Testing Method} is evaluated against it.

\section{Downlink}

The downlink schedule results separate the two methods most clearly. The main
difference is not only how long the schedule is, but how efficiently service is
distributed across users once multiuser coupling becomes difficult.

\subsection{Payload Completion Results and Comparison}

Downlink payload completion shows the strongest benchmark advantage. The
\emph{Training Only Method} finishes with user block counts \((4,3,3)\), while
the \emph{Training + Testing Method} expands the same task to \((9,3,7)\).
That means the main schedule tail is carried by users 0 and 2.

\subsubsection*{Training Only Method}

\begin{itemize}[leftmargin=1.5em,itemsep=0.2em,topsep=0.3em]
    \item \textbf{User 0}: 4 blocks, total \(n=198\), served \(1000/1000\), with 3 skipped blocks.
    \item \textbf{User 1}: 3 blocks, total \(n=140\), served \(2000/2000\), with no skipped blocks.
    \item \textbf{User 2}: 3 blocks, total \(n=199\), served \(3000/3000\), with no skipped blocks.
\end{itemize}

\subsubsection*{Training + Testing Method}

\begin{itemize}[leftmargin=1.5em,itemsep=0.2em,topsep=0.3em]
    \item \textbf{User 0}: 9 blocks, total \(n=427\), served \(1000/1000\), with 7 skipped blocks.
    \item \textbf{User 1}: 3 blocks, total \(n=164\), served \(2000/2000\), with no skipped blocks.
    \item \textbf{User 2}: 7 blocks, total \(n=512\), served \(3000/3000\), with 3 skipped blocks.
\end{itemize}

Figure~\ref{fig:dl-payload-schedule} makes the same point visually. The
benchmark remains compact, with average chosen sub-blocklengths approximately
\((49.5,46.7,66.3)\). The \emph{Training + Testing Method} pushes users 1 and 2
to heavier average sub-blocklengths, approximately \((47.4,54.7,73.1)\), while
also fragmenting users 0 and 2 across more blocks. So the remaining gap is not
just a small timing shift; it is a materially longer and less synchronized
schedule.

\begin{figure}[H]
\centering
\begin{subfigure}[t]{0.48\textwidth}
    \centering
    \includegraphics[width=\textwidth,height=0.27\textheight,keepaspectratio]{20260706_schedule_dl_payload_training_only.png}
    \caption{Training Only Method}
\end{subfigure}
\hfill
\begin{subfigure}[t]{0.48\textwidth}
    \centering
    \includegraphics[width=\textwidth,height=0.27\textheight,keepaspectratio]{20260706_schedule_dl_payload_training_testing.png}
    \caption{Training + Testing Method}
\end{subfigure}
\caption{Downlink payload-completion schedule comparison. The benchmark keeps the schedule compact, while the \emph{Training + Testing Method} stretches users 0 and 2 over a much longer block tail.}
\label{fig:dl-payload-schedule}
\end{figure}

\section{Uplink}

The uplink results are much closer to the benchmark. The
\emph{Training + Testing Method} largely preserves the same per-user schedule
shape and differs mainly through a small blocklength redistribution.

\subsection{Payload Completion Results and Comparison}

Uplink payload completion keeps the same overall front-loaded structure in both
methods: most service is delivered early, and the last block is used to clear a
small remainder. The benchmark finishes with block counts \((2,2,3)\), while
the \emph{Training + Testing Method} gives \((2,3,3)\). The only structural
difference is the extra block for user 1.

\subsubsection*{Training Only Method}

\begin{itemize}[leftmargin=1.5em,itemsep=0.2em,topsep=0.3em]
    \item \textbf{User 0}: 2 blocks, total \(n=68\), served \(1000/1000\), and cleared in 2 blocks.
    \item \textbf{User 1}: 2 blocks, total \(n=119\), served \(2000/2000\), and cleared in 2 blocks.
    \item \textbf{User 2}: 3 blocks, total \(n=183\), served \(3000/3000\), with a small final clean-up block.
\end{itemize}

\subsubsection*{Training + Testing Method}

\begin{itemize}[leftmargin=1.5em,itemsep=0.2em,topsep=0.3em]
    \item \textbf{User 0}: 2 blocks, total \(n=67\), served \(1000/1000\), and cleared in 2 blocks.
    \item \textbf{User 1}: 3 blocks, total \(n=129\), served \(2000/2000\), with one extra block versus the benchmark.
    \item \textbf{User 2}: 3 blocks, total \(n=182\), served \(3000/3000\), with a small final clean-up block.
\end{itemize}

So the benchmark edge is very localized. Users 0 and 2 are essentially matched,
while user 1 is the only case where the \emph{Training + Testing Method}
extends the schedule. Figure~\ref{fig:ul-payload-schedule} shows that both
methods still keep the same basic policy of a large early allocation followed
by a sharply reduced final block.

\begin{figure}[H]
\centering
\begin{subfigure}[t]{0.48\textwidth}
    \centering
    \includegraphics[width=\textwidth,height=0.27\textheight,keepaspectratio]{20260706_schedule_ul_payload_training_only.png}
    \caption{Training Only Method}
\end{subfigure}
\hfill
\begin{subfigure}[t]{0.48\textwidth}
    \centering
    \includegraphics[width=\textwidth,height=0.27\textheight,keepaspectratio]{20260706_schedule_ul_payload_training_testing.png}
    \caption{Training + Testing Method}
\end{subfigure}
\caption{Uplink payload-completion schedule comparison. The two methods follow almost the same structure, with the only material change being one extra block for user 1 under the \emph{Training + Testing Method}.}
\label{fig:ul-payload-schedule}
\end{figure}

\section{Schedule Takeaways}

The scheduling results are most naturally read benchmark-first. In downlink
payload completion, the main shortfall of the \emph{Training + Testing Method}
is schedule fragmentation, especially on users 0 and 2.

In the uplink, the schedule gap is already small. Payload completion differs
mainly through one extra user-1 block. So the present scheduling gap is mostly
a downlink issue, while the uplink is already close to the benchmark at the
per-user schedule level.

\end{document}
